\def\baselinestretch{1}

\chapter{Introduction}\label{ch:introduction}

\def\baselinestretch{1.66}

The history of modern electronics is inseparable from the history of lithography. Since the independent invention of the integrated circuit (IC) by Jack Kilby at Texas Instruments and Robert Noyce at Fairchild Semiconductor in 1958, Moore's Law has governed the relentless progression of the semiconductor industry: the density of components on a chip doubles approximately every two years \cite{Schmidt2012_UltraPrecision}. This exponential trend has translated into dramatic reductions in cost alongside equally dramatic increases in performance. The price of one gigabyte of NAND flash memory, for instance, fell from over \$1000 at the start of this millennium to below \$0.30 by the mid-2010s \cite{Butler2011_PositionControl}, while the minimum transistor feature size has shrunk from hundreds of nanometres to just a few tens of nanometres. The engine of this revolution is \textit{photolithography} — the process by which the intricate patterns of a transistor, a wire, or a memory cell are projected onto a silicon wafer with nanometre-scale fidelity.

The fabrication of a modern IC consists of up to 30 sequential process cycles, each defining one distinct layer of the device \cite{Schmidt2012_UltraPrecision}. A monocrystalline silicon wafer — a thin disc of up to 300\,mm in diameter — serves as the substrate throughout this process. Each cycle begins with a chemical pre-treatment, such as the thermal growth of a thin oxide film for electrical isolation. A photosensitive material called \textit{resist} is then deposited onto the wafer surface. An optical exposure system projects the pattern for the current layer — defined in a chromium film on a transparent quartz plate known as the \textit{reticle}, or \textit{photo-mask} — onto the resist. After development, the regions that received light are selectively dissolved (in the case of a positive resist), exposing the underlying material. This material then undergoes one of several process steps. \textit{Plasma etching} selectively removes oxide or semiconductor material to define trenches and isolation regions. \textit{Ion implantation} locally modifies the electrical properties of the silicon to form transistor junctions and doped wells. \textit{Physical or chemical vapour deposition} lays down thin films of conductive metal — typically tungsten or copper — to form the interconnecting wiring that links the millions of functional elements stacked across all layers. After processing, the remaining resist is stripped and the cycle repeats for the next layer, until all layers are complete and the wafer is diced and packaged.

In the earliest days of the industry, the reticle was brought into direct physical contact with the resist-coated wafer during exposure — a technique known as \textit{contact lithography}. While simple, this approach had two fundamental shortcomings. First, repeated physical contact progressively contaminated and damaged the delicate reticle pattern, necessitating its costly and time-consuming replacement. Second, a 1:1 projection offers no optical resolution advantage: any defect present on the reticle is transferred at full size onto the wafer. These limitations motivated a transition to \textit{projection optical lithography}, in which a reduction lens images the reticle pattern onto the wafer from a safe working distance. The projection lens introduces a demagnification factor of four, meaning that features on the reticle are printed four times smaller on the wafer, immediately relaxing the demands on reticle fabrication. The minimum printable feature size, known as the \textit{critical dimension} (CD), is governed by \cite{Schmidt2012_UltraPrecision}:
\begin{equation}
    \mathrm{CD} = k\,\frac{\lambda}{\mathrm{NA}},
    \label{eq:rayleigh}
\end{equation}
where $\lambda$ is the illumination wavelength, $\mathrm{NA}$ is the numerical aperture of the projection lens, and $k$ is a process-dependent factor. Reducing CD requires a shorter wavelength, a higher NA, and process improvements that lower $k$. This relation has driven a continuous push toward shorter wavelengths — from mercury arc-lamp sources at 436\,nm and 365\,nm, through excimer lasers at 248\,nm and 193\,nm, to the extreme ultraviolet (EUV) sources at 13.5\,nm now entering production.

Achieving a high numerical aperture, however, comes with a practical constraint. Since $\mathrm{NA} = n\sin\vartheta$, where $\vartheta$ is the half-angle of the marginal ray accepted by the lens, attaining values of NA close to one in air requires the lens to subtend a very wide solid angle. Constructing a lens element that simultaneously covers the full area of a 300\,mm wafer at high NA would require an optical element of impractical size and cost \cite{Schmidt2012_UltraPrecision}. The industry's solution was to accept a smaller \textit{exposure field} — typically $26 \times 32$\,mm, corresponding to a single IC die — and expose the wafer in successive steps: after each die is exposed, the wafer stage repositions the wafer to the next die location and the exposure is repeated. This architecture defines the \textit{wafer stepper}. In a stepper, each die is illuminated in a single, instantaneous full-field flash, with the stage brought to a complete stop before exposure commences.

While the stepper represented a decisive step forward, exposing an entire die simultaneously subjects the full field to any non-uniformities in the illumination intensity or lens aberrations. The \textit{scanner} addresses this by replacing the full-field flash with a \textit{scanning exposure}: the illumination system reveals only a narrow rectangular slit of the reticle, and the reticle stage sweeps the mask through this slit while the wafer stage moves synchronously in the opposite direction, maintaining the 4:1 demagnification relationship throughout \cite{Butler2011_PositionControl, Schmidt2012_UltraPrecision}. Because the reticle moves at four times the speed of the wafer, it traverses the full reticle pattern in the same time that the wafer moves one die height. This scanning motion offers significant advantages over full-field exposure: it averages out illumination dose non-uniformities across the exposure field, improves imaging uniformity at the field edges, and allows the use of a physically larger reticle that is easier to write with electron-beam pattern generators. After each die is scanned, the wafer is stepped to the next die position and the scan repeats, defining the \textit{step-and-scan} cycle that gives modern lithographic tools their name. A throughput of over 175 wafers per hour — corresponding to less than 0.2\,s of exposure time per die — is not uncommon in state-of-the-art tools \cite{Butler2011_PositionControl}.

A critical consequence of the multi-layer nature of IC fabrication is the requirement that each new layer be positioned with nanometre-level accuracy relative to all preceding layers — a specification known as \textit{overlay}. The functional layers of a chip are interconnected by millions of tiny conductive pillars called \textit{vias}; a misregistration between consecutive layers of even a few nanometres can misalign a via, break an intended connection, or create an unintended short circuit, rendering the IC defective \cite{Schmidt2012_UltraPrecision}. Overlay requirements scale roughly in proportion with the CD and currently sit below 4\,nm for leading-edge nodes. Overlay is therefore one of the primary technology drivers in scanner design, and it is governed directly by the positioning accuracy of the wafer and reticle stages during the scanning exposure.

The commercial performance of a scanner is ultimately measured in wafers per hour (WPH), since a machine whose capital cost can exceed \$150 million must generate revenue continuously. Although the exposure time per die is fixed by the optics and the process, the time spent stepping and \textit{settling} — the interval between the end of a step motion and the moment the stage has damped its residual vibrations to within the nanometre-level positioning specification — represents a significant and controllable fraction of the total cycle time \cite{Butler2011_PositionControl}, which places the allowable settle time below 10\,ms. Every millisecond saved in settling directly translates into more dies exposed per hour and thus more chips delivered per unit time, with a lower cost per chip. Reducing settling time is therefore a central objective in scanner control engineering.

It is worth noting that, while EUV lithography at 13.5\,nm has entered high-volume production for leading-edge nodes below 10\,nm, deep ultraviolet (DUV) optical lithography at 193\,nm remains the dominant technology for \textit{mature nodes} — technology generations at 28\,nm and above. Mature-node chips are the backbone of automotive electronics, industrial controllers, power management devices, analogue circuits, and the vast majority of consumer electronics. The aggregate volume of wafers produced at mature nodes substantially exceeds that at leading-edge nodes, and is projected to remain so for the foreseeable future. Throughput improvements in DUV step-and-scan tools therefore continue to have a large global impact.

The key lever for reducing settling time — and thus for increasing scanner throughput — is the generation of \textit{reference trajectories}: motion profiles that command the stage from its current position to the next scanning start-point with minimal excitation of mechanical resonances and the shortest possible residual settling tail. A well-designed trajectory limits the energy injected into the structural eigenmodes of the machine, reduces induced vibrations in the sensitive metrology and optical systems, and enables the stage to reach nanometre-level positioning accuracy well within the allocated window. The design of such trajectories is a problem of direct industrial relevance, and forms the central subject of this work.

\medskip

%%% ----------------------------------------------------------------------
